\section*{RESUMEN}

La presente investigación tuvo como objetivo determinar la eficacia de un modelo híbrido basado en Stacking, XGBoost y LightGBM para mejorar la detección de ataques minoritarios en redes del Internet de las Cosas (IoT). El estudio abordó el problema del desbalance de clases en sistemas de detección de intrusiones, debido a que los ataques de baja frecuencia suelen ser clasificados incorrectamente por modelos tradicionales, incrementando el riesgo de falsos negativos. La investigación fue de tipo aplicada, con enfoque cuantitativo, nivel explicativo-comparativo y diseño experimental en entorno controlado. Se utilizó el dataset CICIoT2023, dividido en entrenamiento, validación y prueba, con 5\,491\,971, 1\,176\,851 y 1\,176\,851 registros respectivamente, cada uno con 47 variables. El procedimiento siguió la metodología CRISP-DM e incluyó limpieza de datos, tratamiento del desbalance mediante Borderline-SMOTE, entrenamiento de modelos base XGBoost y LightGBM, generación de probabilidades de clase, cálculo de entropía de Shannon y clasificación final mediante un Árbol de Decisión como meta-modelo. La evaluación se planteó mediante Recall, F1-Score, Accuracy, tiempo de entrenamiento, latencia de inferencia y consumo de memoria RAM, comparando el modelo híbrido con Random Forest y XGBoost como modelos base. Asimismo, se consideró la prueba de rangos con signo de Wilcoxon para validar diferencias significativas. Se espera que la arquitectura propuesta incremente la sensibilidad frente a ataques minoritarios y mantenga una eficiencia computacional adecuada para escenarios de análisis de tráfico IoT.

\textbf{Palabras clave:} Borderline-SMOTE, CICIoT2023, desbalance de clases, detección de intrusiones, entropía de Shannon, IoT, LightGBM, Stacking, XGBoost.

\section*{ABSTRACT}

This research aimed to determine the effectiveness of a hybrid model based on Stacking, XGBoost, and LightGBM to improve the detection of minority attacks in Internet of Things (IoT) networks. The study addressed the class imbalance problem in intrusion detection systems, since low-frequency attacks are often misclassified by traditional models, increasing the risk of false negatives. The research was applied, quantitative, explanatory-comparative, and experimental under a controlled environment. The CICIoT2023 dataset was used, divided into training, validation, and testing subsets with 5,491,971, 1,176,851, and 1,176,851 records respectively, each containing 47 variables. The procedure followed the CRISP-DM methodology and included data cleaning, class imbalance treatment using Borderline-SMOTE, training of XGBoost and LightGBM base models, class probability generation, Shannon entropy calculation, and final classification through a Decision Tree as the meta-model. The evaluation was designed using Recall, F1-Score, Accuracy, training time, inference latency, and RAM consumption, comparing the hybrid model against Random Forest and XGBoost as baseline models. In addition, the Wilcoxon signed-rank test was considered to validate significant differences. The proposed architecture is expected to increase sensitivity to minority attacks while maintaining suitable computational efficiency for IoT traffic analysis scenarios.

\textbf{Keywords:} Borderline-SMOTE, CICIoT2023, class imbalance, intrusion detection, Shannon entropy, IoT, LightGBM, Stacking, XGBoost.
